% Generated draft: requires mathematical and source review.
\documentclass[11pt]{article}
\usepackage[margin=25mm]{geometry}
\usepackage{fontspec}
\setmainfont{Latin Modern Roman}
\usepackage{amsmath,amssymb,mathtools,unicode-math}
\setmathfont{Latin Modern Math}
\usepackage{enumitem,needspace,xurl,hyperref}
\setlength{\parindent}{0pt}
\setlength{\parskip}{5pt}
\setlength{\emergencystretch}{3em}
\newcommand{\N}{\mathbb{N}}
\newcommand{\Z}{\mathbb{Z}}
\newcommand{\Q}{\mathbb{Q}}
\newcommand{\R}{\mathbb{R}}
\newcommand{\C}{\mathbb{C}}
\newcommand{\F}{\mathbb{F}}
\newcommand{\A}{\mathbb{A}}
\newcommand{\PP}{\mathbb{P}}
\newcommand{\E}{\mathbb{E}}
\newcommand{\eps}{\varepsilon}
\newcommand{\dd}{\,\mathrm{d}}
\begin{document}
\section*{UnsolvedMath ID: 700790}
\textbf{Source problem code:} catalogue entry
\subsection{Definitions and mathematical statement}
For a finite group \(G\), let
\[
\operatorname{Cl}(G)=\bigl\{\{gxg^{-1}:g\in G\}:x\in G\bigr\},\qquad k(G)=|\operatorname{Cl}(G)|.
\]
Thus \(k(G)\) counts all conjugacy classes of elements, including the identity class. Write \(|G|\) for the order of \(G\).

The source uses the expression ``finite (simple) group.'' It does not specify whether the questions concern all finite groups, finite simple groups, or both classes separately. Let \(\mathcal C\) denote the intended class; this choice remains unresolved. If simple groups are intended, the source also does not specify whether simplicity is restricted to nonabelian groups.

\begin{enumerate}
\item[(a)] (Attributed in the source to R. Brauer.) Determine the best possible upper bound of the form
\[
\forall G\in\mathcal C,\qquad |G|\leq f\bigl(k(G)\bigr),
\]
where \(f:\Z_{\geq 1}\to\R_{>0}\). The source does not specify whether ``best possible'' means a pointwise optimal bound, an optimal asymptotic estimate, or optimization within a particular family of functions. These interpretations must be distinguished in a definitive formulation.
\item[(b)] The archived question asks whether
\[
\forall G\in\mathcal C,\qquad |G|\leq \exp\bigl(k(G)\bigr),
\]
with the usual convention \(\exp(t)=e^t\).
\end{enumerate}

The supplied Archive of Solved Problems entry reports a negative answer to (b), citing the group \(M_{22}\), with reported order \(443520\) and reported number of conjugacy classes \(12\). It attributes this report to T. Plunkett, \emph{Letter of May, 9, 2000}. This is a supplied archival assessment, not an independently verified result here. The supplied registry labels the overall record partially solved, with archived solved part(s) and remaining part(s) in the main section.
\subsection{Short English statement}
Part (a), attributed to R. Brauer, asks for the sharpest bound on a finite group's order in terms of its number of conjugacy classes. The parenthetical restriction to simple groups and the meaning of sharpness require clarification. Part (b) asks specifically for the bound \(|G|\leq e^{k(G)}\); the supplied archive reports that \(M_{22}\) disproves it.
\subsection{Sources}
\begin{itemize}
\item Local UnsolvedMath registry, permanent ID \(700790\), exact source problem code \texttt{catalogue entry}, title \emph{Kourovka Notebook Problem 11.98}. The supplied record identifies S. P. Strunkov as proposer and gives the year 1990; part (a) is separately attributed to R. Brauer.
\item \emph{The Kourovka Notebook (Unsolved Problems in Group Theory)}, 11th issue (1990), Problem 11.98; the supplied background identifies an update in the 21st issue (2026), arXiv:1401.0300v48. The supplied collection description names E. I. Khukhro and V. D. Mazurov as editors. These bibliographic identifications are supplied source information; the referenced text has not been independently consulted here.
\item Supplied \emph{Archive of Solved Problems} entry for 11.98: negative answer to part (b), reporting that \(M_{22}\) has order \(443520\) and \(12\) conjugacy classes; attributed to T. Plunkett, \emph{Letter of May, 9, 2000}. No separate publication details or link for the letter are supplied.
\item Supplied status review dated \texttt{2026-10-16}: canonical status \texttt{partially\_solved}, evidence label \texttt{published}; its notes place solved part(s) in the archive and remaining part(s) in the main section. This assessment is supplied dated curation, not independently verified.
\end{itemize}
\end{document}
